Nullity of $X(x)$ and orthogonality of $\mathbb{W}$ to $\mathbb{V}$ give
\[
Q_{\mathbb{W}}(f(x)) = -Q(x) \qquad \text{for all } x \in \mathbb{V}.
\]
Setting $x = 0$: $Q_{\mathbb{W}}(f(0)) = 0$ with $f(0) \neq 0$, so $\mathbb{W}$ contains a nonzero null vector. A nondegenerate quadratic space containing a nonzero null vector is indefinite \cite{Lam2005} (a null vector has a hyperbolic partner, \claimref{std-algebra.witt-decomposition}), hence contains directions of both signs and has dimension at least $2$; increments $(+k,0)$ and $(0,+k)$, whose forms are definite, are excluded for every $k$. (When $\mathbb{V}$ has mixed signature the same conclusion follows a second way, with no appeal to the origin: $-Q$ attains both signs, hence so does $Q_{\mathbb{W}}$.)

Sufficiency of one hyperbolic plane is the construction of Sections \ref{sec:model}--\ref{sec:totality}. For uniqueness of the minimal $\mathbb{W}$: a two-dimensional nondegenerate quadratic space containing a nonzero null vector is indefinite, hence of signature $(1,1)$, and is the hyperbolic plane, unique up to isometry (Section \ref{subsec:null}, \cite{Lam2005}; \claimref{std-algebra.sylvester-inertia}). Hence $\dim \mathbb{W} = 2$ with signature $(1,1)$ is minimal and unique, and the increment is $(+1,+1)$.
